\section{Additional Experimental Protocols}
\label{app:experiments}

\begin{figure}[t]
  \centering
  \begin{tikzpicture}[
  inode/.style={infra, font=\small, text width=32mm, minimum height=10mm, inner sep=3pt},
  anode/.style={attack, font=\small, text width=32mm, minimum height=10mm, inner sep=3pt},
  hnode/.style={honest, font=\small, text width=32mm, minimum height=10mm, inner sep=3pt},
  dnode/.style={defense, font=\small, text width=32mm, minimum height=10mm, inner sep=3pt},
  wnode/.style={infra, font=\small, text width=74mm, minimum height=10.5mm, inner sep=3pt}
]
  % Row 1: Datasets & Attack Vectors (y = 0)
  \node[inode] (data) at (0, 0) {\textbf{Datasets \& Tasks}\\[0.3mm]\scriptsize StrategyQA, AESLC, PersonaChat};
  \node[anode] (attack) at (4.4, 0) {\textbf{Attack Vectors}\\[0.3mm]\scriptsize AgentPoison \& MemoryGraft};

  % Row 2: Ingress Channels & Defenses (y = -1.9) -> Gap = 9mm!
  \node[hnode] (channels) at (0, -1.9) {\textbf{6 Mobile Channels}\\[0.3mm]\scriptsize Typed, SMS, notify, share, etc.};
  \node[dnode] (methods) at (4.4, -1.9) {\textbf{Evaluated Defenses}\\[0.3mm]\scriptsize Baselines \& \PG{} (Ours)};

  % Row 3: Physical Hardware & Metrics (y = -3.8) -> Gap = 9mm!
  \node[wnode] (metrics) at (2.2, -3.8) {\textbf{Physical Xiaomi 13T Hardware Measurements}\\[0.3mm]\scriptsize Metrics: ASR, Utility, FQR, Step Latency, Peak RSS, Bit-Exact Recovery};

  % Long, prominent, beautifully visible arrows
  \draw[dataflow] (data.east) -- (attack.west);
  \draw[attackflow] (attack.south) -- (methods.north);
  \draw[dataflow] (data.south) -- (channels.north);
  \draw[dataflow] (channels.east) -- (methods.west);
  \draw[controlflow] (methods.south) -- (methods.south |- metrics.north);
  \draw[dataflow] (channels.south) -- (channels.south |- metrics.north);
\end{tikzpicture}

  \caption{Experimental setup from channel payloads to metrics. Each attack is replayed through each feasible channel while holding text constant, enabling provenance-policy effects to be separated from content detection.}
  \label{fig:experimental-setup}
\end{figure}

\subsection{Factorial design}

A full factorial is expensive, so experiments proceed in stages. Stage A fixes one model/device and crosses six channels with four attacks and two poison budgets. Stage B carries the strongest two attacks across devices, models, and LoRA ranks. Stage C evaluates recovery intervals and replay horizons. Stage D evaluates adaptive retyping with a controlled human or scripted confirmation study. Each stage freezes its analysis plan before inspecting final outcomes.

\begin{table}[h]
\centering
\small
\caption{Planned experimental factors. Values are protocol settings, not results.}
\resizebox{\linewidth}{!}{%
\begin{tabular}{lll}
\toprule
Factor & Primary levels & Secondary levels \\
\midrule
Channel & message, notification, calendar & share, clipboard, screenshot \\
Model & Qwen2.5-0.5B, Gemma-3-1B-class & framework-permitted compact model \\
LoRA rank & 4, 8, 16 & module subset \\
Poison budget & count and percentage & burst versus slow drip \\
Checkpoint interval & 1, 3, 7 windows & adaptive interval \\
Replay horizon & short, medium, full study & encrypted capsule off/on \\
Audit & lineage, behavior, update & combined \\
\bottomrule
\end{tabular}%
}
\end{table}

\subsection{Cross-sink laundering validation}

The killer experiment has four mandatory checkpoints: (1) the payload is present only in an untrusted-channel event; (2) a memory descendant is retrieved; (3) the descendant or its content enters a logged training window; and (4) after source and memory deletion, the active adapter retains trigger-conditioned behavior. A trial lacking the lineage evidence for step (3) is a memory-persistence trial, not a laundering trial.

To rule out cached-context artifacts, restart the process, clear inference caches, disable memory retrieval, and reload the active adapter from disk before the post-deletion measurement. To rule out trigger collision, evaluate the same trigger on the clean base and clean-personalized branches. To establish causal dependence, repeat with the candidate window removed and with a clean replay from the pre-poison checkpoint.

\subsection{Rollback validation}

Compare four states: clean counterfactual $S_C$; attacked $S_A$; approximate-unlearned $S_U$; and PoisonGuard-recovered $S_R$. Check serialized adapter identity where determinism permits, numerical distances $\|\theta_R-\theta_C\|_2$ and $\|\theta_U-\theta_C\|_2$, memory-set equality, fixed-probe logits, ASR, and utility. Exact restoration is accepted only if the predefined state criterion and all lineage completeness checks pass.

\subsection{Systems measurement controls}

Condition batteries to a fixed range, disable unrelated background jobs, randomize method order, allow cool-down, and record ambient temperature. For energy, prefer an external inline monitor; correlate but do not substitute software counters without calibration. Repeat cold- and warm-cache runs. Include failed and thermally throttled trials in an appendix and justify exclusions.

\subsection{Statistical analysis}

Use user-level or conversation-level resampling. Report confidence intervals and the number of independent training seeds. For paired branches, analyze paired differences. Correct for multiple comparisons within each research question or designate a single primary contrast. Avoid declaring ``no utility loss'' from a nonsignificant test; use equivalence margins chosen before measurement.

\subsection{Result integrity checklist}

Before replacing placeholders, verify that every plotted point maps to raw runs; every baseline used its recommended tuning budget; attacks are not tuned on evaluation probes; device and model failures are reported; confidence intervals reflect independent units; and the caption distinguishes measured, derived, and conceptual values. Do not copy performance numbers from cited papers into PoisonGuard result tables.
